\section{Resistance as a Stabilizer}

\subsection{The Physical Anchor}

Computation can advance abruptly while the systems it is meant to transform remain largely recognizable. A factory, a city, and a health service all absorb change through structures that already have work to do. The gap between a new capability and a transformed world is produced inside those structures.

Two processes matter: execution and confirmation. Execution meets resources, process time, and dependencies. Confirmation meets evidence, standards, and responsibility. Together they shape the transition from a technical possibility to something on which people act.

A damped oscillator gives the simplest picture of how a system's response depends on more than the force applied:

\begin{equation}
m\frac{d^2x}{dt^2}+\gamma\frac{dx}{dt}+kx=F(t),
\qquad
\zeta=\frac{\gamma}{2\sqrt{mk}},
\qquad
\tau_{\mathrm{slow}}\approx\frac{\gamma}{k}\quad(\zeta\gg1).
\end{equation}

For a mechanical displacement $x$, $m$ is mass, $\gamma$ the damping coefficient, $k$ the restoring stiffness, and $F(t)$ the driving force. Their units are kg, kg/s, kg/s$^2$, and N respectively. The damping ratio $\zeta$ is dimensionless. In the strongly overdamped regime, the slow relaxation time is approximately $\gamma/k$: more damping produces a longer transition.

As an analogy for technological change, inertia represents installed assets, damping represents resistance to transition, and restoring forces represent arrangements that maintain the existing state. For a constant input, the equilibrium displacement depends on $F/k$; the approach to it depends on the dynamic parameters. A force and a timetable are different parts of the model. Social and technical systems need their own measured variables before this analogy can yield a forecast.

Maintenance makes the distinction concrete. A body needs metabolism; a server needs power and cooling; a transport network needs inspection and repair. A building can persist for a time without an active supply, while its materials continue to age. These systems sustain order through different mechanisms and on different timescales. Work on an operating railway can be harder to schedule than work on an empty site because the railway must continue to carry people.

\subhead{Resource barriers}

Transforming matter requires a physical route. Ironmaking must reduce iron oxides. Cement manufacture must drive reactions and process material at high temperature. Furnaces, feedstocks, equipment, and heat flows belong in the account. Better chemistry and process design can lower actual energy use and sometimes alter the thermodynamic requirement by changing the route. An activation barrier, a reaction's free-energy difference, and a plant's energy consumption are distinct quantities. Computation earns its value by finding better routes through these conditions.

This is the resource barrier, $R$. A bridge requires materials and a site. A grid-control algorithm acts on a network whose generation, transmission, storage, and operating margins already exist. A robot policy runs through a battery, actuators, sensors, joints, and a maintenance system. Intelligence changes how resources are used. The resulting operation still needs resources.

Human bodies belong in the same comparison. Converting a daily food intake of 2,500 dietary kilocalories into average power gives

\begin{equation}
P_{\mathrm{human}}
=\frac{E_{\mathrm{food}}}{\Delta t}
=\frac{2500\times4.18\times10^3}{86400}
\approx1.2\times10^2\ \mathrm{W}.
\end{equation}

That is about 120~W of average metabolic input over a day. It supports the whole organism, including temperature regulation, cognition, and activity. Mechanical output is a different quantity.

A comparison with a machine requires the same task, duty cycle, quality, and system boundary. If a machine consumed an average of 1--3~kW under those conditions, its input-power ratio to 120~W would be about 8--25. Those machine values are an illustration. Equipment, labor, energy prices, downtime, and service life then determine the economic comparison.

A flexible human worker can be highly economical in an irregular setting. In a standardized operation, a specialized machine may reverse the comparison. The relevant protection comes from the present cost of completing the task.

\subhead{Temporal rigidity}

Every transition needs time to unfold. Heat must move, reactions must proceed, material damage must become observable, organisms must respond, and people must coordinate. This is temporal rigidity, $T$.

A candidate molecule and evidence about its long-term effects are separated by exposures and observations. A simulated traffic flow and a functioning transport service are separated by vehicles, maintenance, users, operating permissions, and practice. An infrastructure asset may be designed for decades of use, so a replacement must fit the service it is replacing.\cite{miatto2017how}

Parallel construction, prefabrication, better testing, and better coordination can shorten parts of a project. The useful analysis identifies which part of the critical path actually changes. Compress a year of design work to a day while leaving a multiyear build unchanged, and the number of proposed projects can grow much faster than the number completed.

Waiting also has more than one function. A redundant approval can be removed. An observation period may reveal a delayed effect. Consultation can expose a cost that a project sponsor would otherwise impose on someone else. Calling all three “delay” conceals the work each is doing.

\subhead{Topological coupling}

A system is a network of dependencies. A local improvement must pass through equipment, people, standards, routines, interests, and feedback. This is topological coupling, $G$.

A medical model enters a hospital's diagnostic workflow, record systems, staffing, reimbursement arrangements, and responsibility for treatment. A new energy device enters a grid. A better driving policy enters a road system used by other people. A strong component can fail to transform the whole because the surrounding connections do not support the transition.

Moravec's paradox draws attention to the difficulty of ordinary embodied performance. A formally specified game can demand immense search while still supplying explicit rules and a clear account of success. A robot handling laundry must cope with uncertain shape, occlusion, contact, friction, and deformation. The environment with the more complicated rules on paper can be easier to formalize.

Discretizing $n$ degrees of freedom into $m$ values each produces

\begin{equation}
N_{\mathrm{config}}=m^n,
\qquad
\ln N_{\mathrm{config}}=n\ln m.
\end{equation}

This count explains why exhaustive search rapidly becomes impractical. Geometry, learning, feedback, and passive mechanics can avoid that enumeration. The remaining difficulty depends on the task and the method. A person folding a shirt benefits from a body and a lifetime of interaction already fitted to the problem.

Standardizing the environment changes the comparison. A warehouse with regular containers, predictable shelving, and controlled access can be much easier to automate than an arbitrary home. Part of the achievement has been purchased by rebuilding the task's surroundings. Resource barriers and topological coupling move together.

\subhead{Three media}

Bits, atoms, and life absorb intelligence differently.

In the domain of \term{bits}, generated output can often become visible immediately. Text, images, code, data, and virtual scenes are easy to duplicate and revise within a recoverable environment. Trust imposes another threshold: reliable code, traceable knowledge, and decisions with an accountable owner require more than the arrival of an output.

In the domain of \term{atoms}, a prediction must become a force, a contact, a movement, or a flow. Manufacturing, construction, transport, energy, logistics, and robotics pay for mass, friction, inertia, strength, space, and maintenance. A wrongly cast component needs physical rework. An improved controller changes what machinery can do within those conditions.

In the domain of \term{life}, intervention changes a system that continues to respond. Cells, bodies, organisms, and ecosystems compensate, adapt, and sometimes evolve. AI can accelerate target identification, structural prediction, molecular design, and screening. Toxicology, pharmacokinetics, clinical effects, and sustained safety require evidence from the relevant biological processes. A gene's function depends on its regulatory, developmental, cellular, and ecological setting.

The basic pattern is uneven absorption. Bits support rapid diffusion, atoms require work, and life returns feedback. A model capability can move quickly through one layer and slowly through another.

\begin{figure}[ht]
\centering
\begin{tabularx}{\textwidth}{@{}lXXX@{}}
\toprule
 & \textbf{Bits} & \textbf{Atoms} & \textbf{Life}\\
\midrule
Absorption & Copy and diffuse & Apply force and move material & Respond and adapt\\
Main demands & Compute, evidence, trust & Energy, materials, maintenance & Regulation, exposure, sustained feedback\\
Reversal & Often cheap within a saved state & Repair and reconstruction & Recovery may be limited or delayed\\
Example & A patch in a test branch & A robot handling an object & A drug acting in a body\\
\bottomrule
\end{tabularx}
\caption{Three media, three ways of absorbing an intervention. The same project can span all three columns.}
\label{fig:media-gradient}
\end{figure}


Code acquires a much higher cost of failure when it controls a production system. A drug candidate enters a different timetable when it reaches a trial. The task and the consequence matter alongside the medium.

The three coordinates can trade against one another. More resources may buy time; more time may allow a lower-resource process; redundancy may improve reliability while adding equipment and maintenance. New processes can also reduce several costs together. A lighter user experience can conceal a heavier backend, while a genuine process improvement can reduce the total.

\key{Which resistance fell, and where did the rest go?} That question follows a technology beyond its demonstration.

\subsection{The Grounds for Relying on a Result}

Execution asks how a proposal reaches the world. Confirmation asks what warrants relying on it. A well-formed statement, a supported claim, and a successful intervention require different evidence.

\subhead{The verification pyramid}

A thousand drug candidates do not create a thousand trial programs. A large library of driving policies does not establish how any one behaves on public roads. A polished policy proposal does not establish its effects on organizations or public trust. Verification consumes resources, time, and access to the system being investigated.

Six layers make the progression visible. \term{Formal checks} test syntax, rules, and specified logical properties. \term{Semantic checks} examine meaning, facts, citations, and sources. \term{Causal tests} examine whether an intervention produces an effect. \term{Engineering validation} examines operation under noise, failure, maintenance, cost, and scale. \term{Long-term biological and human evaluation} examines effects that need time to emerge. \term{Accountability} asks whether a deployed arrangement can be explained, audited, challenged, and made to answer for harm.

For candidates filtered through successive stages,

\begin{equation}
N_k=N_0\prod_{i=1}^{k}p_i,
\qquad 0\leq p_i\leq1,
\end{equation}

where $p_i$ is the fraction passing stage~$i$ among those that passed all earlier stages. The expression follows from conditional counting, without an independence assumption. If a stage leaves no survivors, later counts remain zero. It counts survivors; the cost and reliability of each test require separate measures.

\begin{figure}[ht]
\centering
\begin{tikzpicture}[x=1cm,y=1cm,font=\scriptsize]
\foreach \y/\w/\shade/\tag in {0/10.8/12/Formal checks,0.64/9.7/17/Semantic checks,1.28/8.6/22/Causal tests,1.92/7.5/27/Engineering validation,2.56/6.4/32/Long-term evaluation,3.2/5.3/37/Accountability}{
  \fill[deepblue!\shade] (-\w/2,\y) rectangle (\w/2,\y+0.58);
  \node at (0,\y+0.29){\tag};
}
\draw[-{Stealth[length=2mm]},slate] (-5.9,0)--(-5.9,3.8);
\node[rotate=90,slate] at (-6.3,1.9){More context and longer feedback};
\end{tikzpicture}
\caption{Different questions require different evidence. The shape organizes the verification tasks; its widths are illustrative.}
\label{fig:verification-pyramid}
\end{figure}


Tools can improve every layer, at different rates. Automatic checks may increase throughput sharply. Field exposure, long-term observation, and the resolution of disputed consequences remain tied to other processes. When candidate generation outruns them, the waiting room fills.

The upper layers also change the standard of judgment. The acceptance conditions for a formal specification can be explicit. An accountable deployment must respond to the actual distribution of risks and interests. Operational rules can encode part of that task. Treating the encoded portion as the entire objective creates a new failure mode: participants learn to optimize the measure while the purpose recedes.

\subhead{Proof and its premises}

Gödel's incompleteness theorems establish limits for consistent, effectively axiomatized formal systems strong enough to express basic arithmetic. Under the relevant hypotheses, some propositions cannot be decided within the system; the second theorem restricts its ability to prove its own consistency.\cite{godel1931}

An engineering example brings the practical question closer. A proof that a program satisfies a specification leaves open whether the specification captures the intended requirement. Machines can check proofs and help check specifications. Sensors, tests, and independent evidence can extend the chain. The remaining judgment concerns what errors the chain covers and what its omissions would cost.

For decisions that allocate risk and benefits, standards also need authorization. People who bear the consequences need ways to provide evidence, contest a requirement, grant permission, and appeal a decision. Formal competence and legitimate authority answer different questions.

\subhead{A promise needs an address}

Trust includes a question that an accuracy score leaves unanswered: where does the loss go when the decision fails?

Nassim Taleb's idea of \emph{skin in the game} draws attention to the relation between decision makers and those exposed to their decisions. A person can pledge property, professional standing, and reputation. An organization can supply assets, a durable identity, insurance, and a structure for correction and compensation. These arrangements give a promise an address.

Medical and other consequential services need both reliable decisions and a workable response to failure. Records, review, suspension authority, and a responsible operator help establish that response. A human assigned to approve outputs needs the information, time, and power to intervene. A name on a form supplies little by itself.

Trust can also accumulate through a relationship. Care, counseling, and companionship involve memory, knowledge of a situation, and commitments that are kept over time. A model can generate a response. The service still has to establish continuity, crisis handling, and correction.

Evidence gives a judgment grounds; accountability gives a commitment a return path. Both can be supported by automation and institutions. Together with the conditions for execution, they stretch the transition from a new capability to a changed society.

The transition is a contested interval. Work can be reorganized on one timetable while people's account of their own worth changes on another. That mismatch is the next problem.
